\chapter[Christliche Heilwissenschaft]{Die sogenannte „Christliche Heilwissenschaft“}
\quellseite{7}
(Christian Science), die mit ihren schriftwidrigen und unsinnigen Lehren schon so viele Christen bethört hat und noch immer neue Anhänger gewinnt, ist kürzlich von einer indischen Philosophin, Pundita Ramabai, die America besuchte, ganz richtig beurtheilt worden. Bekanntlich lehren die Anhänger dieses Irrwahns, daß nichts Körperliches wirklich existire, daß darum auch Krankheit und Schmerzen des Leibes nur Einbildungen seien. Man kämpfe nur gegen diese Einbildungen und überwinde sie, dann sei man gesund. Dazu bemerkt nun die genannte Indierin, daß dies nichts anderes als alte heidnische, indische Philosophie sei, und sagt: „Ich bin verwundert und erstaunt, daß bei allem Fortschritt des 19. Jahrhunderts alte Philosophien im Gewande christlicher Namen in den Vereinigten Staaten auftreten. Es ist ein trauriger Anblick für einen, der mit den Folgen heidnischer Philosophie und heidnischen Aberglaubens bekannt ist, zu sehen, wie gebildete Leute, die alle Vorrechte christlicher Civilisation genießen, durch den Glanz eines neuen Namens betrogen werden. Als ich im Frühjahr in New York landete, hörte ich, daß in America eine neue Philosophie gelehrt würde und schon viele Anhänger habe. Man nannte sie ‘Christian Science’, und als ich mich nach ihren Lehren erkundigte, fand ich, daß es dieselbe Philosophie war, die unter meinem Volk schon 4000 Jahre gelehrt wird. Da ich darin geboren und unterrichtet bin, bin ich mit ihren Schriften und ihrem Einfluß auf das Volk bekannt und möchte ihre Verderbtheit bezeugen.“ Pundita Ramabai gab dann als Inhalt dieser Philosophie dies an: „Du sollst das ganze Weltall für nichts als Betrug ansehen. Du mußt glauben, es sei gar nicht vorhanden. Du selbst existirst nicht. Ich existire nicht. Wenn du das erfaßt hast, hast du die Weisheit.“ Von den Americanern sagte sie: „Ihr seid ein Volk, das Mitgefühl hat. Alles ist wirklich. Ihr fühlt, daß ihr, wenn andere verhungern wollen, ihnen zu essen geben müßt. In Indien ist das ganz anders. Da hat niemand Mitleid mit dem andern, da ist kein Mitgefühl für die Verhungernden oder Verwundeten. In unserer letzten Hungersnoth hatten unsere Philosophen kein Erbarmen mit den Leidenden und halfen niemand. Warum sollten sie ihnen auch helfen, wenn sie lehren, alles Leiden sei Einbildung und der Tod der Kinder sei Einbildung? Die erste Folge dieser falschen Weisheit ist also die verwerflichste Grausamkeit: kein Mitleid mit den Leidenden und höchste Selbstsucht.“ — Und doch gibt es auch unter lutherischen Christen solche, die diesen „kräftigen Irrthümern und Lügen“, 2 Thess. 2, 11., glauben, gerade wie auch der widerlichste aller Lügenpropheten, der marktschreierische Dowie, Leute aus unsern Gemeinden zum Abfall verführt hat. So groß ist die Macht des Irrthums und seines Urhebers, des Teufels. L. F.
